\documentclass[11pt,a4paper]{article}
\usepackage[utf8]{inputenc}
\usepackage[T1]{fontenc}
\usepackage{lmodern}
\usepackage[french]{babel}
 
\usepackage[top=2.2cm, bottom=2.2cm, left=2.2cm, right=2.2cm]{geometry}
\usepackage{fancyhdr}
\usepackage{microtype}
\usepackage{setspace}
\setstretch{1.15}

 
\usepackage{amsmath, amsfonts, amssymb}
\usepackage{siunitx}
\sisetup{locale = FR, detect-all}

 
\usepackage{graphicx}
\usepackage{booktabs}
\usepackage{tabularx}
\usepackage{xcolor}
\usepackage{float}
\usepackage{tikz}
\usetikzlibrary{shapes.geometric, arrows.meta, positioning, calc}

 
\definecolor{primary}{RGB}{0,0,0}      % Bleu marine académique
\definecolor{secondary}{RGB}{0,0,0}  % Bleu cyan technique
\definecolor{darkgray}{RGB}{50, 50, 50}
\definecolor{lightgray}{RGB}{245, 247, 250}

 
\usepackage[colorlinks=true, linkcolor=primary, citecolor=secondary, urlcolor=secondary]{hyperref}
 
\begin{document}

 
\begin{center}
    {\LARGE \textbf{\textcolor{primary}{Modélisation d'un générateur PV avec Commande MPPT pilotée par un Réseau de Neurones.}}}
 
 
    \rule{\linewidth}{1.2pt}
\end{center}

\vspace{0.3cm}

 

\vspace{0.5cm}

 
 
\section{Modélisation Mathématique du Panneau PV}
Le panneau est modélisé à l'aide du modèle standard à une diode, représenté par l'équation suivant:
\begin{equation}
    I_{\mathrm{pv}} = I_{\mathrm{ph}} - I_{\mathrm{rs}} \left[ \exp\left( \frac{q \, V_{\mathrm{pv}}}{p \, k \, T_k \, n_s} \right) - 1 \right]
    \label{eq:pv_diode}
\end{equation}
où $q = \SI{1.602e-19}{\coulomb}$ est la charge élémentaire, $k = \SI{1.38e-23}{\joule\per\kelvin}$ la constante de Boltzmann, $n_s = 36$ le nombre de cellules en série, $p = 1.3$ le facteur d'idéalité de la jonction, et $T_k$ la température absolue en Kelvin.

Le photocourant $I_{\mathrm{ph}}$ et le courant de saturation inverse $I_{\mathrm{rs}}$ sont asservis aux conditions environnementales réelles :
\begin{align}
    I_{\mathrm{ph}} &= \left[ I_{\mathrm{sc}} + \alpha (T_k - T_{\mathrm{ref}}) \right] \cdot \frac{G}{1000} \\
    I_{\mathrm{rs}} &= I_r \left( \frac{T_k}{T_{\mathrm{ref}}} \right)^3 \exp\left[ \frac{q \, E_g}{p \, k} \left( \frac{1}{T_{\mathrm{ref}}} - \frac{1}{T_k} \right) \right]
\end{align}

\paragraph{Solveur Numérique:} 
La méthode de Newton-Raphson est implémentés pour calculer le courant $I_{pv}$ à chaque étapes de temps. Pour chaque tension $V_pv(idx)$, $I_{pv}$ est obtenue par itération:
\begin{equation}
    I^{(k+1)} = I^{(k)} - \frac{f(I^{(k)})}{f'(I^{(k)})}, \quad \text{avec } f'(I) = 1
\end{equation} 
Le Solveur stop lorsque $|\Delta I| > 10^{-8}$, les courbes \textit{I - V} et \textit{P - V} nous permet de localiser le maximum power point:
\begin{equation}
    P_{\mathrm{max}} = \max(V_{\mathrm{pv}} \times I_{\mathrm{pv}}) \implies (V_{\mathrm{mpp}}, I_{\mathrm{mpp}})
\end{equation}
Pour simplifier le calcul numérique on prend $R_{shunt} = 0$
\section{Convertisseur DC-DC Boost}

Pour transférer la puissance maximale vers une charge résistive $R_L$, un hacheur élévateur (Boost) est inséré entre le générateur PV et la charge.

En régime de conduction continue (CCM), la relation de transfert de tension idéale est :
\begin{equation}
    V_{\mathrm{load}} = \frac{V_{\mathrm{pv}}}{1 - D}
\end{equation}
En négligeant les pertes du convertisseur ($P_{\mathrm{in}} \approx P_{\mathrm{out}}$), l'impédance d'entrée équivalente $R_{\mathrm{in}}$ vue par le panneau solaire s'exprime en fonction de la résistance de charge $R_L$ et du rapport cyclique $D \in [0, 1[$ :
\begin{equation}
    R_{\mathrm{in}} = (1 - D)^2 R_L
    \label{eq:impedance}
\end{equation}
Pour contraindre le panneau à opérer exactement à son point de puissance maximale, l'impédance d'entrée doit égaler l'impédance interne optimale du panneau $R_{\mathrm{in,mpp}} = \frac{V_{\mathrm{mpp}}}{I_{\mathrm{mpp}}}$. On en déduit analytiquement la loi de commande optimale :
\begin{equation}
    D_{\mathrm{mpp}} = 1 - \sqrt{\frac{R_{\mathrm{in,mpp}}}{R_L}} = 1 - \sqrt{\frac{V_{\mathrm{mpp}}}{I_{\mathrm{mpp}} \cdot R_L}}
    \label{eq:duty_cycle}
\end{equation}
Cette relation analytique fournit la consigne idéale de rapport cyclique servant de vérité terrain pour l'entraînement du réseau neuronal.
  
\section{Réseau de Neurones (FFNN)}

\subsection{Constitution de la Base de Données d'Apprentissage}
Le modéle mathématique développés sous MATLAB fait un balayage couvrant l'ensemble du domaine d'exploitation:

\begin{itemize}
    \item Ensoleillement $G \in [100, 1000]\,\si{\watt\per\meter\squared}$ avec un pas de $\SI{50}{\watt\per\meter\squared}$.
    \item Température $T \in [15, 65]\,\si{\celsius}$ avec un pas de $\SI{5}{\celsius}$.
\end{itemize}
Pour chaque pair (G,T), la puissance maximale et le rapport cyclique sont calculés via l'équation \ref{eq:duty_cycle}. L'ensemble de données généré comporte 209 échantillons.
\subsection{Architecture et Topologie du Réseau}
Le réseau choisi est un Perceptron Multicouche (MLP) à propagation avant (\textit{Feedforward}) :
\begin{center}
\begin{tikzpicture}[x=1.5cm, y=0.9cm, >=Latex]
    % Couche d'entrée
    \node[circle, draw=primary, fill=primary!15, line width=1pt, minimum size=0.8cm] (in1) at (0, 1) {$G$};
    \node[circle, draw=primary, fill=primary!15, line width=1pt, minimum size=0.8cm] (in2) at (0, -1) {$T$};
    \node[above=0.2cm of in1, font=\bfseries\small, text=primary] {Entrées};

    % Couche cachée
    \foreach \i in {1,...,5} {
        \node[circle, draw=secondary, fill=secondary!20, line width=1pt, minimum size=0.6cm] (h\i) at (2.5, 2.5 - \i*0.8) {};
    }
    \node at (2.5, -2.1) {$\vdots$};
    \node[circle, draw=secondary, fill=secondary!20, line width=1pt, minimum size=0.6cm] (h10) at (2.5, -2.9) {};
    \node[above=0.2cm of h1, font=\bfseries\small, text=secondary] {Couche Cachée (10 neurones)};
    \node[below=0.1cm of h10, font=\scriptsize\itshape] {Activation : tansig};

    % Couche de sortie
    \node[circle, draw=primary, fill=primary!15, line width=1pt, minimum size=0.8cm] (out) at (5, 0) {$D$};
    \node[above=0.2cm of out, font=\bfseries\small, text=primary] {Sortie};
    \node[below=0.1cm of out, font=\scriptsize\itshape] {Activation : purelin};

    % Connexions
    \foreach \i in {1,2,3,4,5,10} {
        \draw[->, gray!70] (in1) -- (h\i);
        \draw[->, gray!70] (in2) -- (h\i);
        \draw[->, secondary] (h\i) -- (out);
    }
\end{tikzpicture}
\end{center}

\begin{itemize}
    \item \textbf{Couche d'entrée :} Irradiance $G$ et Température $T$.
    \item \textbf{Couche cachée  :}  fonction d'activation tangente sigmoide.
    \item \textbf{Couche de sortie:} Rapport cyclique $D \in [0.05, 0.85]$ avec fonction linéaire ($\mathrm{purelin}$).
\end{itemize}

\subsection{Optimisation et Validation}
L'apprentissage est fait via l'algorithme de \textbf{Levenberg-Marquardt} (trainlm) (trés bon choix pour les problèmes d'approximation de fonction).
Les données sont reparties comme ceci: 70\% apprentissage, 15\% validation, 15\% test indépendant.


\section{Simulink}

Le réseau entraîné a été exporté sous Simulink (\texttt{gensim}) et interconnecté dans le modèle de puissance complet multiphysique (\texttt{circuit\_pv.slx}) comprenant le panneau solaire, l'inductance de lissage ($L = \SI{0.4}{\milli\henry}$), la capacité intermédiaire ($C = \SI{50}{\micro\farad}$), l'interrupteur MOSFET modulé à une fréquence de découpage $F_{\mathrm{sw}} = \SI{10}{\kilo\hertz}$, et la charge résistive.

Un profil d'ensoleillement variable en échelons ($500 \to 1000 \to 800 \to 600 \to \SI{200}{\watt\per\meter\squared}$) a été appliqué pour confronter le contrôleur neuronal au contrôleur conventionnel à incrément de conductance (InCond).
 \\

\noindent\textbf{Analyse des résultats :}
\begin{enumerate}
    \item \textbf{Rapidité dynamique :} Le réseau de neurones calcules instantanément par produit matricielle la valeur optimale de $D$ 
    \item \textbf{Robustesse   :} Un test a été menée: On applique un ensoleillement étrange sur notre système ($G > 1000$ $ W \text{m}^{-2}$) on constate que le réseau de neurones gére avec aisance l'Irradiance inconnue (gràce à la fonction tangente sigmoide) et le systéme reste stable en toutes cas. 

\end{enumerate}
 
\section{Conclusion et Compétences Développées}

Ce projet valide l'efficacité de l'association entre intelligence artificielle et électronique de puissance pour la commande de systèmes énergétiques renouvelables. 
\end{document}
